\section{Arguments from computational complexity and computability}
\label{sec:complexity}

A second family of arguments asks whether the laws of physics as we know them
could be run on a computer at all, and at what cost. Complexity arguments ask
how resources scale. Computability arguments ask whether any algorithm can do
the job. Before reviewing the literature we separate three tasks that are often
conflated:
\begin{enumerate}
\item[(T1)] \emph{Execution}: evolving a given initial state forward by the
  dynamical law, to some accuracy, for a finite time.
\item[(T2)] \emph{Decision}: answering a question about a model that quantifies
  over infinitely many times, sizes or instances, such as whether the system
  ever halts, whether a spectral gap survives the thermodynamic limit, or
  whether a ground state has energy below $E_0$.
\item[(T3)] \emph{Proof}: deriving every true sentence of a formal theory of
  the model from its axioms.
\end{enumerate}
A simulator of a universe must perform (T1). Most of the negative results below
concern (T2) or (T3). A negative result for (T2) or (T3) bears on the
simulation hypothesis only through an additional premise linking it to (T1). We
state that premise explicitly in each case.

\subsection{Feynman, Lloyd and the extended Church-Turing thesis}
\label{sec:complexity:feynman}

Feynman's keynote address concerned \emph{exact} simulation, in which ``the
computer will do exactly the same as nature''. He explicitly set aside
approximate numerical integration~\cite{Feynman1982simulating}. He imposed the
efficiency rule that the number of computer elements should be proportional to
the simulated space-time volume. He was candid that this was a stipulation:
``I make up the rules, I'm allowed to do that''~\cite[p.~469]{Feynman1982simulating}.
Under that rule, a classical machine that stores the general probability
function of $N$ variables fails, because doubling $N$ requires an
``exponentially explosive growth'' of the computer~\cite[p.~472]{Feynman1982simulating}.
His categorical ``No!'' concerned a different question. He asked whether a
\emph{local} classical probabilistic computer can reproduce quantum
probabilities. The obstruction he identified is a Bell-type
inequality~\cite[pp.~476, 485]{Feynman1982simulating}. He conjectured that
locally coupled quantum spins could imitate any discrete quantum system. He
was confident for bosons but left fermions open~\cite[p.~476]{Feynman1982simulating}.

Lloyd showed that the conjecture holds for local Hamiltonians. A quantum
computer simulates a closed local system with memory and time ``directly
proportional to the time and space taken up by the system to be
simulated''~\cite{Lloyd1996universal}. Continuous variables, such as a lattice
gauge field, are handled by discretization, with $N$ qubits per variable giving
$N$ bits of accuracy~\cite{Lloyd1996universal}. Bravyi and Kitaev showed that
fermionic modes and qubits simulate each other with at most logarithmic
overhead per gate~\cite{Bravyi2002fermionic}. Polynomial-time quantum
algorithms exist for scattering in $\phi^4$ theory in up to four spacetime
dimensions~\cite{Jordan2012quantum}. Lloyd regarded such results as ``a strong
indication'' that any eventual theory of quantum gravity will be efficiently
simulable on a quantum computer~\cite{Lloyd2002computational}. That is an
expectation, not a theorem.

The complexity-theoretic premise behind ``the universe is too hard to
simulate'' is the failure of the extended Church-Turing thesis (ECT). The ECT
holds that any reasonable model of computation can be simulated efficiently by
a probabilistic Turing machine. Unlike the original thesis, it is an empirical
hypothesis whose truth depends on physics~\cite{Copeland2023churchturing}.
Bernstein and Vazirani gave the first formal evidence against it: an oracle
problem solvable in quantum polynomial time but requiring superpolynomial time
classically. They also showed that $\mathsf{BPP}\subseteq\mathsf{BQP}\subseteq
\mathsf{P}^{\#\mathsf{P}}$. An unrelativized proof of quantum advantage would
therefore require a major breakthrough in complexity
theory~\cite{Bernstein1997quantum}. Experimental evidence is likewise
conditional. Arute et al.\ estimated that a sampling task their processor
performed in about 200\,s would take a state-of-the-art classical supercomputer
about 10{,}000 years~\cite{Arute2019quantum}. A later tensor-network method
produced one million samples for that circuit, at fidelity $F\approx0.0037$, in
about 15 hours on 512 GPUs~\cite{Pan2022solving}.

What follows validly from these results? If the ECT fails in our universe, a
\emph{classical} simulator that faithfully reproduces quantum dynamics incurs a
superpolynomial overhead on some instances. A quantum simulator does not, by
Lloyd's theorem. The argument therefore constrains the simulator's architecture
and running time. It does not bear on whether a simulation is possible.

\subsection{The sign problem}
\label{sec:complexity:sign}

Troyer and Wiese defined a ``solution'' of the sign problem as a
polynomial-complexity algorithm for a thermal average in a model whose related
bosonic problem is itself polynomial. A change of representation that leaves
the cost exponential does not count~\cite{Troyer2005computational}. They mapped
the NP-complete ground-state problem of the three-dimensional $\pm J$ Ising
spin glass to a quantum model with a sign problem,
$H=-\sum J_{jk}\sigma^x_j\sigma^x_k$. The related bosonic model is a
ferromagnet that cluster algorithms simulate efficiently. A \emph{generic}
solution of the sign problem would therefore solve an NP-complete problem in
randomized polynomial time. The paper states P$\,=\,$NP, and a footnote
correctly refines this to BPP~\cite{Troyer2005computational}. The result is
worst-case. It ``does not exclude that a specific sign problem can be solved
for a restricted subclass'', and meron-cluster algorithms are an
example~\cite{Troyer2005computational}.

Two further points in the paper matter here. First, the hard instances require
inverse temperatures growing with $N$. Classical Monte Carlo on the
corresponding spin glass has autocorrelation times $\tau\propto e^{aN}$.
Troyer and Wiese also judge that quantum simulators are ``most likely not a
generic solution'' either~\cite{Troyer2005computational}. Second, physical
relaxation towards minimum-energy states, as in spin glasses, soap films and
protein folding, is itself subject to local optima and long relaxation
times~\cite{Aaronson2005np}. NP-hardness of the sign problem is therefore a
statement about the worst-case cost of a (T2)-type equilibrium task. It does
not show that nature efficiently computes something that a computer cannot. A
simulator that executes the dynamics (T1) need not solve the spin-glass
ground-state problem, just as a physical spin glass does not.

Ringel and Kovrizhin proved a narrower and stronger statement for a specific
algorithm class~\cite{Ringel2017quantized}. Take bosonic systems whose gapped
phases have a nonvanishing quantized thermal Hall conductance, as most
fractional quantum Hall phases do. For these, no \emph{local} sign-free QMC
formulation exists. Here ``local'' means that the configuration variables are
obtained from the physical degrees of freedom by finite-depth local quantum
circuits. The argument extends to a chiral Kagome antiferromagnet under
``additional microscopic assumptions''. The authors state that they ``do not
address the possibility of nonlocal approaches to QMC, such as determinant QMC
or cluster algorithms''. They call obstruction to classical simulation in
general ``a rather ill-defined task''. They concede that, for non-symmetric
transfer matrices, a needed similarity transformation might be nonlocal,
``rendering the previous arguments void''~\cite{Ringel2017quantized}.
Related rigorous results include Hastings' intrinsic sign problem for commuting
double-semion Hamiltonians, under a technical assumption~\cite{Hastings2016how}.
Later work gave criteria excluding stoquastic Hamiltonians and sign-free
determinantal QMC for broad classes of chiral topological
phases~\cite{Smith2020intrinsic,Golan2020intrinsic}. Marvian et al.\ showed
that even \emph{finding} a curing transformation within restricted
single-qubit classes is NP-complete~\cite{Marvian2019on}. None of these
results concerns classical algorithms outside the QMC family, or quantum
simulators.

The Science Advances article does not mention the simulation hypothesis; we
searched the full text for this. Press coverage nevertheless ran headlines
such as ``Physicists Confirm That We're Not Living In a Computer Simulation'',
asserting that it is ``impossible to model the physics of our universe on even
the biggest computer''~\cite{Eck2017physicists}. One report attributed to the
paper quotations about memory ``built from more atoms than there are in the
universe'' that do not occur in the published text~\cite{UPI2017were}.
Aaronson identified four independent gaps between the paper and the headline.
Only local transformations within QMC are excluded. $\mathsf{BPP}\neq\mathsf{BQP}$
is unproven. A simulator could be quantum. An exponential classical slowdown
would be invisible to the simulated~\cite{Aaronson2017because}. Bostrom made
the quantum-computer point in the same context~\cite{Bostrom2025simulation}.

\subsection{Undecidability of physical properties}
\label{sec:complexity:undecidability}

Cubitt, Perez-Garcia and Wolf constructed translationally invariant,
nearest-neighbour Hamiltonians on a 2D lattice of $d$-level systems, with
interaction strength at most~1. Each Hamiltonian is gapped with $\gamma\ge1$
if a universal Turing machine halts on input $n$, and gapless otherwise. The
spectral-gap problem is therefore undecidable~\cite{Cubitt2015undecidability,Cubitt2022undecidability}.
It follows that for any consistent, recursively axiomatized theory there is
such a Hamiltonian whose gap can be neither proved nor
disproved~\cite{Cubitt2022undecidability}. The authors emphasize that
``algorithmic undecidability always concerns infinite families of systems''.
The property at stake is defined in the thermodynamic limit. Their models can
look gapless at all sizes up to a threshold that is itself
uncomputable~\cite{Cubitt2015undecidability}. A recent review by the same group
and co-authors puts the general point plainly. Undecidability ``is a feature
of the mathematical model'', and because any finite problem is decidable, an
undecidable one ``must necessarily conceal an infinite somewhere''. Its finite
``reflections'' take the form of computational-complexity hardness or
size-dependent effects at uncomputably large
sizes~\cite{PeralesEceiza2025undecidability}.

Earlier results have the same structure. Moore embedded a Turing machine in
the motion of a single particle with three degrees of freedom. For such a
system ``virtually any question about [its] long-term dynamics is
undecidable'' even for exactly known initial conditions. Yet finite-time
quantities remain computable, and ``the best one can do is to simulate the
system and see what happens''~\cite{Moore1990unpredictability}. Wolfram put
the general point directly. The behaviour of a physical system ``may always be
calculated by simulating explicitly each step''. Questions about idealized limits
of infinite time, volume or precision can be formally undecidable. The finite versions are decidable, though PSPACE-complete
for universal cellular automata~\cite{Wolfram1985undecidability}. Pour-El and
Richards exhibited computable initial data for a 3D wave equation whose (weak)
solution is not computable~\cite{PourEl1981wave}. We rely here on the account
in Ref.~\cite{PeralesEceiza2025undecidability}. Weihrauch and Zhong showed
that the wave propagator is computable on $C^1$ data (losing one derivative)
and on Sobolev spaces. They argue that the Pour-El-Richards example probably
does not yield a wave computer that beats the Turing
machine~\cite{Weihrauch2002is}. A different kind of limit concerns
\emph{inference devices} embedded in a universe. Wolpert proved that every
such device fails to infer some function, even if it has super-Turing power,
and that no two distinguishable devices can infer each
other~\cite{Wolpert2008physical}. These results limit prediction and
observation from inside a universe. On our reading they say nothing about
whether its dynamics could be executed by an external machine.

What follows validly is that undecidability of a property (T2) is compatible
with, and in the cited constructions is \emph{derived from}, computable local
dynamics (T1). These proofs work by embedding a universal Turing machine in a
computable model~\cite{PeralesEceiza2025undecidability}. An argument from
undecidability to non-simulability therefore needs a further premise: that
executing a single step of the physical dynamics requires the answer to an
undecidable question. That premise amounts to hypercomputation in nature.

\subsection{The Faizal-Krauss-Shabir-Marino argument}
\label{sec:complexity:faizal}

Faizal et al.\ argue from incompleteness to the impossibility of
simulation~\cite{Faizal2025consequences}. Their premises, in order, are as
follows.
(i)~A candidate quantum-gravity theory is a computational formal system
$\mathcal F_{\rm QG}$. Its axioms are recursively enumerable; on the same page
they are also said to be finite. It is arithmetically expressive, consistent
and ``empirically complete''.
(ii)~Gödel's first theorem gives $\mathrm{Th}(\mathcal F_{\rm QG})\subsetneq
\mathrm{True}(\mathcal F_{\rm QG})$, with truth defined as satisfaction in the
standard model of arithmetic. Tarski's and Chaitin's theorems add further
limits.
(iii)~``Physically these Gödel sentences correspond to empirically meaningful
facts—e.g., specific black-hole microstates.''
(iv)~Reality instantiates an external, non-recursively-enumerable truth
predicate $T(x)$ that is ``actualized in nature''. Undecidability results such
as the spectral gap supply ``empirical backing'', and the Lucas-Penrose
argument supplies motivation.
(v)~Simulation proposals assume ``that every physical truth is reducible to the
output of a finite algorithm''.
They conclude that ``any finite algorithm can at best emulate'' the computable
fragment, that reality ``cannot be instantiated on a Turing-equivalent
device'', and that the simulation hypothesis is ``logically impossible rather
than merely implausible''~\cite{Faizal2025consequences}. A university press
release described this as showing the idea to be ``mathematically impossible''
and as possibly ``the final, definitive answer''~\cite{UBCO2025physicists}.

The valid core is conditional. \emph{If} the observable dynamics of our
universe depends on a non-recursively-enumerable predicate, then no
Turing-equivalent machine reproduces it. The antecedent is a denial of the
physical Church-Turing thesis. Gödel's theorems do not supply it. Step~(ii)
concerns provability of arithmetic sentences within a theory, which is task
(T3). Step~(iii), which carries the argument from (T3) to physics, is asserted
without derivation. Redden's comment makes the corresponding objection. The
argument conflates what a formal theory can prove with what an algorithm can
execute. A Turing-complete cellular automaton such as Conway's Game of Life
has undecidable long-term questions yet evolves by a finite
rule~\cite{Redden2025provability}. Non-simulability would require dynamics
that must resolve undecidable problems in order to advance, i.e.\
hypercomputation, for which there is no physical evidence~\cite{Redden2025provability}.
We add three points. First, the physical undecidability results offered as
``empirical backing'' are theorems about infinite families of computable
models, proved by embedding Turing machines in them
(\S\ref{sec:complexity:undecidability}). They exhibit algorithmic rules
generating undecidable global properties, not non-algorithmic dynamics.
Second, even granting all premises, the conclusion excludes only
Turing-equivalent simulators. A host universe with hypercomputational
resources is not addressed, so ``logically impossible'' overstates the paper's
own result. Third, premise~(v) is not Bostrom's position: he holds that only
what simulated observers could notice need be computed~\cite{Bostrom2003are}.
In a 2026 reply to a published discussion paper, the authors clarified that
they take Gödelian truths to be ``objectively actualized in physical
reality'', independently of human cognition~\cite{Faizal2026reply}. This
restates the contested premise; it does not argue for it. A formal
computer-science treatment reaches the opposite verdict under the physical
Church-Turing thesis. Using Kleene's second recursion theorem, Wolpert proved
that a universe satisfying the thesis and its converse can simulate itself.
He also showed that many questions about simulation are undecidable by Rice's
theorem~\cite{Wolpert2025what}. The disagreement therefore reduces to whether
physical dynamics is Turing-computable. That is an open empirical question,
not a consequence of incompleteness.

\subsection{Hypercomputation and the physical Church-Turing thesis}
\label{sec:complexity:hypercomputation}

The Church-Turing thesis concerns what a human can compute by rote. The
physical theses of Deutsch and Wolfram, that every finitely realizable
physical system can be simulated by a universal machine, are distinct
empirical claims~\cite{Copeland2023churchturing}. Idealized Newtonian machines
with unbounded speeds violate Gandy's ``Thesis M''. Deutsch argued that a
universal Turing machine cannot perfectly simulate continuous classical
dynamics~\cite{Copeland2023churchturing}. Aaronson argued that Planck-scale
physics and the holographic entropy bound appear to exclude Zeno-type and
Malament-Hogarth hypercomputers, as well as unlimited-precision analog
computation. He suggested that the intractability of NP-complete problems
might itself be treated as a physical principle~\cite{Aaronson2005np}. For the
present review, the symmetry of the situation is what matters. A
computability-based argument against simulation requires our physics to
violate the physical Church-Turing thesis. If it does, nothing in the argument
prevents the host's physics from doing the same.

\subsection{Resource bounds and the host's physics}
\label{sec:complexity:resources}

Bostrom estimated $\sim\!10^{33}$--$10^{36}$ operations for a realistic
simulation of human mental history, against $\sim\!10^{42}$ operations per
second for a planetary-mass computer of known design. He also held that
simulating the whole universe ``down to the quantum level is obviously
infeasible, unless radically new physics is
discovered''~\cite{Bostrom2003are} (see \S\ref{sec:philosophy:bostrom}). Lloyd
bounded the computation the observable universe can have performed at
${\lesssim}10^{120}$ operations on ${\sim}10^{90}$ bits (${\sim}10^{120}$ with
gravitational degrees of freedom). He read these numbers both as upper bounds
on computation by matter within the universe and as lower bounds for a quantum
computer simulating it directly~\cite{Lloyd2002computational}. Vazza concluded
that simulating our universe, or even a low-resolution Earth, is impossible for
a simulator in a universe ``sharing the same properties''. He added that
``only universes with very different physical properties'' could produce ours
as a simulation, a question he places outside what is scientifically
testable~\cite{Vazza2025astrophysical}.

The logical point is that such bounds bind a simulator only if the simulator
is subject to our physics. Several authors state this explicitly. Neukart et
al.\ derive that a simulating computer ``governed by the same physical laws''
would need at least as many particles as the simulated universe, unless the
programmer ``isn't limited by the simulation's physical
laws''~\cite{Neukart2022do}. Edge and Brown note that Vazza's conclusions hold
``if the simulator's universe is governed by physics like our
own''~\cite{Edge2026commentary}. Bostrom remarks that ``the physics in the
basement universe might allow for vastly more powerful
computers''~\cite{Bostrom2025simulation}. Wolpert's framework allows
simulating and simulated universes to obey different laws~\cite{Wolpert2025what}.
Running time is similarly unconstrained, because a slowdown in the host is not
observable from inside~\cite{Aaronson2017because}. Resource arguments therefore
constrain two scenarios: Bostrom-type ancestor simulations run by our
descendants, whose physics is ours by hypothesis, and nested simulations
within a single physics. They do not constrain the unrestricted hypothesis.
The unrestricted hypothesis, in turn, gains immunity from these arguments only
at the cost of testability (\S\ref{sec:methodology:unrestricted}).

\begin{table*}[t]
\centering
\small
\caption{Complexity and computability arguments: what each establishes about
simulability.}
\label{tab:complexity}
\begin{tabular}{p{0.2\linewidth}p{0.36\linewidth}p{0.36\linewidth}}
\hline
Argument & Establishes & Does not establish \\
\hline
Feynman; Lloyd & Exponential cost for classical storage of general quantum states; efficient quantum simulation of local Hamiltonians & That classical simulation is impossible; the ECT's failure is unproven \\
Troyer-Wiese & No generic polynomial-time sign-problem solution unless NP $\subseteq$ BPP & That nature solves these instances efficiently \\
Ringel-Kovrizhin & No local sign-free QMC for chiral bosonic phases & Anything about nonlocal QMC, other classical algorithms, or quantum simulators \\
Cubitt et al.; Moore; Wolfram & Undecidable thermodynamic-limit or infinite-time properties of computable models & Uncomputability of finite-time dynamics \\
Faizal et al. & Incompleteness of r.e.\ theories (standard) & That physical dynamics is non-algorithmic; exclusion of non-Turing hosts \\
Bostrom; Lloyd; Vazza & Resource limits for simulators within our physics & Limits on hosts with different physics \\
\hline
\end{tabular}
\end{table*}

Table~\ref{tab:complexity} summarizes the section. No complexity-theoretic or
computability result reviewed here rules the simulation hypothesis in or out.
Taken together, and granting their unproven premises (for example,
$\mathsf{BPP}\neq\mathsf{BQP}$), they narrow the admissible designs. A
simulator of quantum physics within our physics would have to be quantum or
pay superpolynomial overhead. It would have to approximate continuum quantities
to finite precision. It would never need to decide the undecidable properties
of the models it runs.
